% ══════════════════════════════════════════════════════════════════════════════
\chapter{Traitement des requêtes}
\label{chap:td5}
% ══════════════════════════════════════════════════════════════════════════════

La cinquième étape constitue la couche d'entrée du système de recherche : il s'agit de
transformer une requête en langage naturel français — saisie librement par un utilisateur
— en une structure de données exploitable par les index inversés. Le cours (chapitre~4,
Grammaires d'interrogation) distingue deux niveaux dans ce traitement : l'extraction des
\textbf{métadonnées structurées} (dates, rubriques, filtres) et l'identification des
\textbf{opérateurs logiques}. L'approche adoptée s'appuie sur des expressions régulières
hiérarchisées, préférées aux analyseurs de dépendances spaCy qui se sont révélés
inefficaces sur ce domaine spécialisé (0/12 entités DATE reconnues par
\texttt{fr\_core\_news\_sm}).

\section{Architecture du parser de requêtes}

La classe \texttt{QueryParser} applique cinq étapes séquentielles. Chacune reçoit le
texte résiduel laissé par la précédente, garantissant qu'aucun span déjà consommé
ne pollue les étapes suivantes. La sortie est un dataclass \texttt{ParsedQuery} figé :

\begin{lstlisting}[caption={Pipeline d'extraction séquentiel de \texttt{QueryParser.parse()}},
  label={lst:query_pipeline}]
def parse(self, query: str) -> ParsedQuery:
    text = query.strip()
    date_min, date_max, text = _extract_dates(text)       # 12 patterns regex
    rubrique, text           = _extract_rubrique(text)    # 6 sections ADIT
    filtre_images, zone, text = _extract_filters(text)    # images / zone titre
    operateurs, text         = _extract_operators(text)   # AND / OR / NOT
    mots_cles                = _extract_keywords(text, self._spell_checker)
    return ParsedQuery(mots_cles=mots_cles, rubrique=rubrique,
                       date_min=date_min, date_max=date_max,
                       operateurs=operateurs, filtre_images=filtre_images,
                       zone=zone)
\end{lstlisting}

L'ordre des étapes est critique : les filtres d'images sont extraits \emph{avant} les
opérateurs pour que le mot \texttt{sans} dans \og{}sans image\fg{} soit consommé par la
règle des filtres et non interprété comme un opérateur NOT.

\section{Extraction des métadonnées et des opérateurs}

\paragraph{Dates.} Douze patterns compilés, ordonnés du plus spécifique au plus général.
Une plage \og{}en 2014\fg{} produit \texttt{date\_min='2014-01-01'} et
\texttt{date\_max='2014-12-31'} ; un mois seul (\og{}novembre 2011\fg{}) calcule le
dernier jour via \texttt{calendar.monthrange}. Toutes les dates sont normalisées en
ISO~8601 pour être directement comparables lors du filtrage en TD6.

\paragraph{Rubriques.} Six sections ADIT reconnues avec leurs alias (variantes accentuées,
singulier/pluriel). Le contexte lexical (\og{}de la rubrique\fg{}, \og{}provenant de\fg{}…)
est consommé pour éviter les résidus dans les mots-clés.

\paragraph{Opérateurs logiques.} Le pattern NOT couvre les marqueurs explicites
(\texttt{mais pas}, \texttt{et non}, \texttt{sans}) et la négation discontinue française
\texttt{ne [verbe] pas} — forme courante dans les requêtes naturelles, invisible pour
un matching lexical simple. OR est détecté par \verb|\bou\b|. En l'absence de tout
marqueur, \texttt{operateurs = ["AND"]} par défaut.

\section{Résultats}

\begin{table}[ht]
\centering
\caption{Exemples de requêtes traitées (extrait des 33 requêtes de l'annexe TD5)}
\label{tab:td5_examples}
\begin{tabular}{>{\small}p{5.8cm} >{\small}p{2.8cm} >{\small}p{2cm} >{\small}p{1.8cm}}
\toprule
\textbf{Requête} & \textbf{mots\_clés} & \textbf{rubrique} & \textbf{opérateurs} \\
\midrule
articles de la rubrique HE parlant des systèmes embarqués
  & [systèmes, embarqués] & HE & AND \\
parus entre le 3 mars 2013 et le 4 mai 2013
  & [] & — & AND \\
Focus, santé, 2014
  & [santé] & Focus & AND \\
CNRS mais pas Centrale Paris
  & [cnrs, centrale, paris] & — & AND, NOT \\
avec des images dont le titre contient croissance
  & [croissance] & — & AND \\
\bottomrule
\end{tabular}
\end{table}

Le tableau~\ref{tab:td5_examples} montre que le pipeline traite correctement 31 des
33 requêtes de l'annexe. Les deux limitations résiduelles sont l'extraction d'une seule
rubrique simultanée (requête~[28]) et l'exclusion du token \texttt{3D} par le tokeniseur
lettre-only. La couverture de tests atteint 96\,\% sur 352 tests (68 spécifiques à ce TD).
